\documentclass[12pt,leqno]{article}
\usepackage{amsfonts,amsmath,amssymb}
\usepackage{color,graphicx}
\usepackage{fullpage,setspace,multirow}
\usepackage[colorlinks=true,urlcolor=darkgray,bookmarks=false]{hyperref}

%\setlength{\parindent}{0mm}
\addtolength{\parskip}{2mm}

\def\changemargin#1#2{\list{}{\rightmargin#2\leftmargin#1}\item[]}
\let\endchangemargin=\endlist

\definecolor{scarlet}{cmyk}{0,1.00,0.65,0.15}
\definecolor{gray}{cmyk}{0.06,0,0,0.34}
\definecolor{darkgray}{cmyk}{0.07,0,0,0.90}

\begin{document}

\vspace*{-20mm}

\begin{changemargin}{-0.5in}{0.5in}
\footnotesize{
\begin{tabular}{lr}
\multirow{1}{*}{\includegraphics[height=0.66in,width=0.66in]{Syllabus/osu_logo.pdf}} & \vspace{-5pt}\textbf{\small{Department of Economics}}\\
&\color{scarlet}{\rule{6.375in}{0.4pt}}\hspace*{1pt}\\
& 410 Arps Hall\\
& 1945 North High Street\\
& Columbus, OH 43210-1172\\
& \\
\textbf{ } & \textbf{Prof. Paul J. Healy}\\
  & Email: \href{mailto:healy.52@osu.edu}{healy.52@osu.edu}\\
\end{tabular}
}
\end{changemargin}


%\bigskip

\begin{center}
{\Large \textbf{Syllabus}\\ECON 8877 - Experimental Economics Methodology}\\\smallskip {\large Autumn 2025}
\end{center}

\[
\begin{tabular}[c]{r|l}
\textbf{Contact Info}&  \\
\hline
\textbf{Name}                           & Prof. Paul J. Healy         \\
\textbf{Office}                         & Arps 465A             \\
\textbf{E-mail (Preferred)}                         & healy.52@osu.edu      \\
\textbf{Phone}                          & 614-247-8876              \\
 \multicolumn{2}{c}{ } \\
 \textbf{Course Info}\\
\hline
\textbf{Course Meeting Times}                       & Tue \& Thur 2:20pm--3:40pm \\
\textbf{Meeting Location} & Arps Hall, Room 318 \\
\textbf{Office Hours}                   & Tue 9:30--10:30am  \\
\textbf{Course Website}                 &\texttt{http://carmen.osu.edu} \\
\textbf{Credit Hours} & 3 \\
\textbf{Mode of Delivery} & in-person
\end{tabular}
\]

\paragraph{Prerequisites:}

This course is designed for second-year students in the Economics PhD program. Students are expected to have taken ECON 8714. Exceptions to this prerequisite can be given at the permission of the instructor.

\paragraph{Course Objective:}

This course will equip students with the tools needed to design, run, and analyze a laboratory experiment that conforms to the norms and expectations of the field of experimental economics. Broadly, the course will cover theoretical discussions of incentives in experiments, including elicitation techniques, as well as empirical research on those methods. Students will also learn relevant statistical techniques for experiments that are not taught in the standard econometrics course sequence. Next, students will receive an introduction to various programming interfaces used to design experiment software. Finally, students will learn and discuss ethical considerations when running experiments, including sample size calculations, $p$-hacking, and the appropriate use of pilot experiments.

\paragraph{Learning Objectives:}

By the end of the course students should be able to take a given research proposal, design an incentive-compatible experiment that is both feasible to run and addresses the proposed research question. They should be able to formulate and execute a data analysis plan. They should know how to calculate the correct sample size for their given hypothesizes effect sizes. Students should know how to submit an IRB approval and apply for funding. They should have the tools needed to begin the process of programming an experiment. And they should know how to handle changes to the design as needed, what to include as pilot data, and what to report in the final paper.

\paragraph{Recommended Texts:}

There are no required texts, but the following are recommended:
\begin{enumerate}
  \item \emph{Experimetrics}, by Peter Moffatt. ISBN-13: 978-0230250222
  \item \emph{Nonparametric Statistics for the Behavioral Sciences (2nd ed.)}, by Sidney Siegel and John Castellan, Jr. ISBN-13: 978-0070573574.
\end{enumerate}

No other materials are required for this course.

\paragraph{Assignments \& Grades:}

There are four types of activities on which you will be assessed:

\begin{description}
\item[1. Data Analysis Exercises (40\%):] Students will be assigned four different data analysis exercises. In each, I will provide a dataset and require students to perform a specific data analysis exercise on that dataset. Examples include maximum likelihood estimation of competing models of preferences, bootstrapping estimates and standard errors, comparing fixed effects vs. random effects, or running a variety of statistical tests. Students may collaborate, but must turn in separate answers. Grading will be based on completeness and accuracy of the analysis. Each assignment will be given equal weight.

\item[2. Problem Set (20\%):] I will assign one problem set based on the decision-theoretic portion of the course. Students will have two weeks to complete the assignment. Students may collaborate, but must turn in separate answers. Grading will be based on correctness of responses.

\item[3. Paper Proposal (20\%):] Each student must submit a draft of an in-progress research project by the midpoint of the semester. Students are encouraged to discuss their ideas with others, but each must generate their own proposal. Students will also give a brief (15--20 minutes) in-class presentation of their proposal. Grading will be based on thoroughness of the proposal and whether the proposed design is feasible and addresses the stated research question.

\item[4. Participation (20\%):] Students are expected to lead the discussion of one topic, and participate in all discussions. Grading will be based on recorded frequency of verbal participation in class. Students are expected to contribute verbally at least twice per week.
\end{description}

Students' letter grades will be determined by the sum of their grades on each of the four components. The final grading scale will be as follows:
\begin{center}
  \begin{tabular}{l|l}
    Grade & Percent Range\\
    \hline    
A	& \ \ 100\%	to	93\% \\
A-	& $<$ 93\%	to	90\% \\
B+	& $<$ 90\%	to	87\% \\
B	& $<$ 87\%	to	83\% \\
B-	& $<$ 83\%	to	80\% \\
C+	& $<$ 80\%	to	77\% \\
C	& $<$ 77\%	to	73\% \\
C-	& $<$ 73\%	to	70\% \\
D+	& $<$ 70\%	to	67\% \\
D	& $<$ 67\%	to	60\% \\
E	& $<$ 60\%	to	0\%
  \end{tabular}
\end{center}

\paragraph{Course Policies:}

Attendance will not be enforced, but strongly encouraged. Students who miss class will likely receive a lower participation grade. Assignments may be turned in via email if needed, but deadline extensions must be approved in advance, except in cases of medical or family emergencies.

\paragraph{Student-Chosen Topics:}

The last few lectures of the course I can adapt to whatever topics students might be interested in learning. Please let me know in advance if there's a topic you'd like to learn about and I will do my best to add it to the schedule.


\paragraph{CarmenCanvas Access:}

All instructional materials will be posted on CarmenCanvas. The login page is at
\href{https://carmen.osu.edu}{https://carmen.osu.edu}. If you need instructions on how to use CarmenCanvas, go to \href{https://teaching.resources.osu.edu/toolsets/carmencanvas/guides/getting-started-carmen-students}{https://teaching.resources.osu.edu/toolsets/carmencanvas/guides/getting-started-carmen-students}.

If you still need help you can contact the IT Service Desk at 614-688-4357 (HELP) and IT support staff will work out a solution with you.

\paragraph{Weather or Other Short-Term Closings:}

Should in-person classes be canceled by the University, we will meet virtually via CarmenZoom during our regularly scheduled time. I will share any updates via email sent through the CarmenCanvas system.


\paragraph{Academic Misconduct:}
Academic integrity is essential to maintaining an environment that fosters excellence in teaching, research, and other educational and scholarly activities. Thus, The Ohio State University and the Committee on Academic Misconduct (COAM) expect that all students have read and understand the University's \href{https://trustees.osu.edu/bylaws-and-rules/code}{Code of Student Conduct}, and that all students will complete all academic and scholarly assignments with fairness and honesty. Students must recognize that failure to follow the rules and guidelines established in the University's Code of Student Conduct and this syllabus may constitute Academic Misconduct.

The Ohio State University’s Code of Student Conduct (Section 3335-23-04) defines academic misconduct as: Any activity that tends to compromise the academic integrity of the University or subvert the educational process. Examples of academic misconduct include (but are not limited to) plagiarism, collusion (unauthorized collaboration), copying the work of another student, and possession of unauthorized materials during an examination. Ignorance of the University’s Code of Student Conduct is never considered an excuse for academic misconduct, so please review the Code of Student Conduct and, specifically, the sections dealing with academic misconduct.

If an instructor suspects that a student has committed academic misconduct in this course, the instructor is obligated by University Rules to report those suspicions to the Committee on Academic Misconduct. If COAM determines that a student violated the University’s Code of Student Conduct (i.e., committed academic misconduct), the sanctions for the misconduct could include a failing grade in the course and suspension or dismissal from the University.

If students have questions about the above policy or what constitutes academic misconduct in this course, they should contact the instructor.

\paragraph{Artificial Intelligence and Academic Integrity:}
There has been a significant increase in the popularity and availability of a variety of generative artificial intelligence (AI) tools, including ChatGPT, Sudowrite, and others. These tools will help shape the future of work, research and technology, but when used in the wrong way, they can stand in conflict with academic integrity at Ohio State.

All students have important obligations under the Code of Student Conduct to complete all academic and scholarly activities with fairness and honesty. Our professional students also have the responsibility to uphold the professional and ethical standards found in their respective academic honor codes. Specifically, students are not to use unauthorized assistance in the laboratory, on field work, in scholarship, or on a course assignment unless such assistance has been authorized specifically by the course instructor. In addition, students are not to submit their work without acknowledging any word-for-word use and/or paraphrasing of writing, ideas or other work that is not your own. These requirements apply to all students undergraduate, graduate, and professional.

To maintain a culture of integrity and respect, these generative AI tools should not be used in the completion of course assignments unless an instructor for a given course specifically authorizes their use. Some instructors may approve of using generative AI tools in the academic setting for specific goals. However, these tools should be used only with the explicit and clear permission of each individual instructor, and then only in the ways allowed by the instructor.

\paragraph{Religious Accommodations:}
Ohio State has had a longstanding practice of making reasonable academic accommodations for students’ religious beliefs and practices in accordance with applicable law. In 2023, Ohio State updated its practice to align with new state legislation. Under this new provision, students must be in early communication with their instructors regarding any known accommodation requests for religious beliefs and practices, providing notice of specific dates for which they request alternative accommodations within 14 days after the first instructional day of the course. Instructors in turn shall not question the sincerity of a student’s religious or spiritual belief system in reviewing such requests and shall keep requests for accommodations confidential.

With sufficient notice, instructors will provide students with reasonable alternative accommodations with regard to examinations and other academic requirements with respect to students’ sincerely held religious beliefs and practices by allowing up to three absences each semester for the student to attend or participate in religious activities. Examples of religious accommodations can include, but are not limited to, rescheduling an exam, altering the time of a student’s presentation, allowing make-up assignments to substitute for missed class work, or flexibility in due dates or research responsibilities. If concerns arise about a requested accommodation, instructors are to consult their tenure initiating unit head for assistance.  

A student’s request for time off shall be provided if the student’s sincerely held religious belief or practice severely affects the student’s ability to take an exam or meet an academic requirement and the student has notified their instructor, in writing during the first 14 days after the course begins, of the date of each absence. Although students are required to provide notice within the first 14 days after a course begins, instructors are strongly encouraged to work with the student to provide a reasonable accommodation if a request is made outside the notice period. A student may not be penalized for an absence approved under this policy.

If students have questions or disputes related to academic accommodations, they should contact their course instructor, and then their department or college office. For questions or to report discrimination or harassment based on religion, individuals should contact the 
\href{mailto:equity@osu.edu}{Civil Rights Compliance Office}.

Policy: \href{https://oaa.osu.edu/religious-holidays-holy-days-and-observances}{Religious Holidays, Holy Days and Observances}

\paragraph{Disability Statement (with Accommodations for Illness):}
The university strives to maintain a healthy and accessible environment to support student learning in and out of the classroom. If students anticipate or experience academic barriers based on a disability (including mental health and medical conditions, whether chronic or temporary), they should let their instructor know immediately so that they can privately discuss options. Students do not need to disclose specific information about a disability to faculty. To establish reasonable accommodations, students may be asked to register with Student Life Disability Services (see below for campus-specific contact information). After registration, students should make arrangements with their instructors as soon as possible to discuss your accommodations so that accommodations may be implemented in a timely fashion.

If students are ill and need to miss class, including if they are staying home and away from others while experiencing symptoms of viral infection or fever, they should let their instructor know immediately. In cases where illness interacts with an underlying medical condition, please consult with Student Life Disability Services to request reasonable accommodations.

\noindent\href{mailto:slds@osu.edu}{slds@osu.edu}\\
\href{https://slds.osu.edu}{https://slds.osu.edu}\\
098 Baker Hall, 113 W. 12th Ave.\\
614-292-3307 phone

\paragraph{Intellectual Diversity:}
Ohio State is committed to fostering a culture of open inquiry and intellectual diversity within the classroom. This course will cover a range of information and may include discussions or debates about controversial issues, beliefs, or policies. Any such discussions and debates are intended to support understanding of the approved curriculum and relevant course objectives rather than promote any specific point of view. Students will be assessed on principles applicable to the field of study and the content covered in the course. Preparing students for citizenship includes helping them develop critical thinking skills that will allow them to reach their own conclusions regarding complex or controversial matters.

\paragraph{Grievances and Solving Problems:}
According to University Policies, if you have a problem with this class, you should seek to resolve the grievance concerning a grade or academic practice by speaking first with the instructor or professor. Then, if necessary, take your case to the department chairperson, college dean or associate dean, and to the provost, in that order. Specific procedures are outlined in Faculty Rule 3335-8-23. Grievances against graduate, research, and teaching assistants should be submitted first to the supervising instructor, then to the chairperson of the assistant’s department.

\paragraph{Creating an Environment Free from Harassment, Discrimination, and Sexual Misconduct:}
The Ohio State University is committed to building and maintaining a welcoming community. All Buckeyes have the right to be free from harassment, discrimination, and sexual misconduct. Ohio State does not discriminate on the basis of age, ancestry, color, disability, ethnicity, gender, gender identity or expression, genetic information, HIV/AIDS status, military status, national origin, pregnancy (childbirth, false pregnancy, termination of pregnancy, or recovery therefrom), race, religion, sex, sexual orientation, or protected veteran status, or any other bases under the law, in its activities, academic programs, admission, and employment. Members of the university community also have the right to be free from all forms of sexual misconduct: sexual harassment, sexual assault, relationship violence, stalking, and sexual exploitation.

To report harassment, discrimination, sexual misconduct, or retaliation and/or seek confidential and non-confidential resources and supportive measures, contact the Civil Rights Compliance Office (CRCO):

Online reporting form: \href{http://civilrights.osu.edu/}{http://civilrights.osu.edu/}

Call 614-247-5838 or TTY 614-688-8605

\href{mailto:civilrights@osu.edu}{civilrights@osu.edu} 

The university is committed to stopping sexual misconduct, preventing its recurrence, eliminating any hostile environment, and remedying its discriminatory effects. All university employees have reporting responsibilities to the Civil Rights Compliance Office to ensure the university can take appropriate action:

\begin{itemize}
    \item All university employees, except those exempted by legal privilege of confidentiality or expressly identified as a confidential reporter, have an obligation to report incidents of sexual assault immediately.
    \item The following employees have an obligation to report all other forms of sexual misconduct as soon as practicable but at most within five workdays of becoming aware of such information: 1. Any human resource professional (HRP); 2. Anyone who supervises faculty, staff, students, or volunteers; 3. Chair/director; and 4. Faculty member.
\end{itemize}

\paragraph{Counseling and Consultation Services / Mental Health:}
As a student you may experience a range of issues that can cause barriers to learning, such as strained relationships, increased anxiety, alcohol/drug problems, feeling down, difficulty concentrating and/or lack of motivation. These mental health concerns or stressful events may lead to diminished academic performance or reduce a student’s ability to participate in daily activities. The Ohio State University offers services to assist you with addressing these and other concerns you may be experiencing.

If you or someone you know are suffering from any of the aforementioned conditions, you can learn more about the broad range of confidential mental health services available on campus via the Office of Student Life’s Counseling and Consultation Service (CCS) by visiting \href{https://ccs.osu.edu}{ccs.osu.edu} or calling 614-292-5766. CCS is located on the 4th floor of the Younkin Success Center and 10th floor of Lincoln Tower. You can reach an on-call counselor when CCS is closed at 614-292-5766 and 24-hour emergency help is also available through the 24/7 by dialing 988 to reach the Suicide and Crisis Lifeline.

\pagebreak

\paragraph{Schedule (Tentative):}
\mbox{}\\
\begin{center}
\begin{tabular}{c|l}
  \textbf{Week} & \textbf{Topics \& Assignments} \\
  \hline\hline
  01 & Theory: Incentives in Experiments I \\ \hline
  01 & Theory: Incentives in Experiments II \\ \hline
  02 & Theory: Model Testing \& Minimal Experiments \\ \hline
  02 & Hypothesis Tests: Fey \& Proschan (2010), $p$-values, units of observation \\ \hline
  03 & Siegal \& Castellan, scales, contingency table tests, correlation tests \\ \hline
  03 & Tests of distributional differences \\ \hline
     & \textbf{Data Analysis Exercise (DAE) 1 assigned} \\ \hline
  04 & Multiple hypothesis testing, Bonferroni-like corrections \\ \hline
  04 & Design \& Ethics: power calculations, $p$-hacking \\ \hline
     & \textbf{DAE 1 due, DAE 2 assigned} \\ \hline
  05 & Belief Elicitation: Scoring rules and incentive compatibility, Savage (1971) \\ \hline
  05 & Belief Elicitation: Multiple price lists, empirical evidence \\ \hline
     & \textbf{DAE 2 due} \\ \hline
  06 & Student presentations I \\ \hline
     & \textbf{Paper Proposals due} \\ \hline
  06 & Student presentations II \\ \hline
  07 & Regressions: Clustering vs Fixed effects vs Random effects \\ \hline
     & \textbf{Problem Set assigned} \\ \hline
  07 & Regressions: Robust standard errors, Gillen Snowberg \& Yariv (2019) \\ \hline
  08 & Regressions: ANOVA vs. dummy variable regressions, interactions, interactions in logit/probit \\ \hline
  08 & Bootstrapping, permutation tests \\ \hline
  09 & Maximum likelihood estimation I \\ \hline
     & \textbf{Problem Set due} \\ \hline
  09 & Maximum likelihood estimation II \\ \hline
     & \textbf{DAE 3 assigned} \\ \hline
  10 & Finite mixture models \\ \hline
  10 & Model selection: BIC, AIC, Cross validation \\ \hline
     & \textbf{DAE 3 due, DAE 4 assigned} \\ \hline
  11 & Writing IRB Approvals \\ \hline
  11 & Getting funded  \\ \hline
     & \textbf{DAE 4 due} \\ \hline
  12 & Programming I: oTree  \\ \hline
  12 & Programming II: zTree \\ \hline
  13 & Programming III: PHP/HTML/JavaScript \\ \hline
  13 & Programming IV: PHP/HTML/JavaScript Part 2 \\ \hline
  14 & Application: Repeated games\\ \hline
  14 & Application: Real-effort tasks \\ \hline
\end{tabular}
\end{center}


\end{document}
